\clearpage
\fancyhead[L]{\small GAUSS AN VON LINDENAU.}
\hypertarget{native-306-lower}{}
\begin{center}GAUSS AN VON LINDENAU.
\end{center}

\begin{center}
\rule{4cm}{0.4pt}\\
\fontsize{10}{12}\selectfont Monatliche Correspondenz. Band XVIII. S. 83..86. Juli 1808.\\
\rule{4cm}{0.4pt}
\end{center}

\begin{flushright}
Göttingen, am 27. Juni 1808.
\end{flushright}

\noindent —\;—\;— Obwohl mein Werk schon im Messekatalog stand, ist der Druck leider bis jetzt noch nicht einmal zur Hälfte abgeschlossen. Es geht sehr langsam voran, da ich mir die einzelnen B\SourceLeaf{307}ögen zur letzten Durchsicht zuschicken lasse, Bei einer Aufgabe, die den Gegenstand eines Abschnitts dieses Werks bildet, nämlich die Bahn eines Planeten aus \emph{vier} geozentrischen Positionen zu bestimmen (von denen zwei unvollständig sein können), habe ich schon vor langer Zeit eine weitere Berechnung der Bahn der \emph{Vesta} angestellt, um ein eigens ausgewähltes Beispiel zu geben. Dabei habe ich eine der letzten Mailänder Beobachtungen mitverwendet. Der Natur der Sache nach konnte und sollte diese Bahn zwar nicht das genauestmögliche Ergebnis aus den vorhandenen Beobachtungen sein (denn alle verwendeten Positionen sind \emph{einzelne} Beobachtungen, die zuvor nicht vom wahrscheinlichsten Beobachtungsfehler befreit wurden). Dennoch war abzusehen, dass diese Elemente viel genauer sein mussten als die III$^{\text{ten}}$, die im September bereits um $1'$ von den Beobachtungen abwichen. Da diese neuen Elemente bisher noch nicht bekannt geworden sind, teile ich sie Ihnen hier mit.


\begin{center}\fontsize{9}{11}\selectfont\setlength{\tabcolsep}{3pt}\renewcommand{\arraystretch}{1.15}\begin{tabular}{lr}
Epoche, Meridian von Paris 1807 .\;.\;.\;.\;.\;.\;.\;.\;.\;.\;.\;.\; & $168^\circ\;10'\;47''6$\\
Tägl. mittl. tropische Bewegung .\;.\;.\;.\;.\;.\;.\;.\;.\;.\;.\;.\;.\;.\; & $978''8588$\\
\shortstack[l]{Aphel für die Epoche\\gegenüber den Fixsternen unverändert} .\;.\;.\;.\;.\;.\;.\;.\; & $69^\circ\;57'\;6''5$\\
Aufsteigender Knoten .\;.\;.\;.\;.\;.\;.\;.\;.\;.\;.\;.\;.\;.\;.\;.\; & $103\;11\;57''3$\\
Exzentrizität .\;.\;.\;.\;.\;.\;.\;.\;.\;.\;.\;.\;.\;.\;.\;.\;.\;.\;.\;.\;.\; & $0{,}0880158$\\
Logarithmus des mittleren Abstands .\;.\;.\;.\;.\;.\;.\;.\;.\;.\;.\;.\; & $0{,}3728980$\\
Neigung der Bahn .\;.\;.\;.\;.\;.\;.\;.\;.\;.\;.\;.\;.\;.\;.\;.\;.\; & $7^\circ\;8'\;20''8$\\
\end{tabular}\end{center}


\par\indent Am 22. Juni haben wir Vesta hier wiedergefunden. Seither hatten wir zwar keine klare Nacht mehr, aber der beobachtete Stern stand fast genau dort, wo Vesta zu erwarten war. Er hatte die achte Größe (wie sie der Planet jetzt ungefähr haben muss), und es gab in der Nähe keinen anderen erkennbaren Stern, der Vesta hätte sein können. Professor \textsc{Harding} hatte bei seiner früheren Durchsicht dieser Gegend hier keinen unbeobachteten Fixstern bemerkt. Schließlich scheinen meine eigenen Beobachtungen über ungefähr eine Stunde das Vorrücken in Rektaszension schon ziemlich unverkennbar zu zeigen. So bleibt wohl kein Zeifel, dass es wirklich Vesta war. Die erste klare Nacht wird uns darüber volle Gewissheit geben. Juno haben wir schon am 20. Juni wiedergefunden und beobachtet, obwohl wir sie an diesem Abend eher in einem anderen Stern 10. Größe zu erkennen glaubten, der aber am 21$^{\text{sten}}$ unverändert an seiner Stelle stand. Am 22$^{\text{sten}}$ gab die Beobachtung volle Gewissheit. Juno erreicht kaum die 10. Größe. Doch hoffe ich, dass Sie sie mit dem Passageinstrument beobachten werden, da sie 1806 bei ungefähr gleicher Lichtschwäche in Mailand mit dem Mauerquadranten beobachtet wurde.

\Needspace{10\baselineskip}
\par\indent Hier sind meine Beobachtungen von Juno und Vesta:


\begin{center}\fontsize{9}{11}\selectfont\setlength{\tabcolsep}{3pt}\renewcommand{\arraystretch}{1.15}\begin{tabular}{lrrrr}
& 1808 & \shortstack{Mittlere Zeit\\in Göttingen} & \shortstack{Scheinbare\\Rektaszension} & \shortstack{Scheinb. südl.\\Deklination}\\\hline
$\JunoSymbol$ & Juni 20 & $11^{\text{u}}\,49^{\text{m}}\,0^{\text{s}}$ & $315^\circ\;29'\;34''$ & $2^\circ\;16'\;23''$\\
$\JunoSymbol$ & 22 & $12\;0\;45$ & $315\;23\;1$ & $2\;14\;20$\\
\VestaAltarThreeOhSeven & 22 & $13\;46\;25$ & $353\;12\;17$ & $9\;11\;55$\\
\end{tabular}\end{center}


\par\indent Bei Juno geben die 7$^{\text{ten}}$ Elemente (Gött. gelehrte Anz. 1808, St. 14, Monatl. Corresp. 1808, Febr.) also die Rektaszension um $11'$ zu klein und die Deklination um $2'$ zu groß an. Bei Vesta geben die dritten Elemente die Rektaszension um $8'$ zu klein und die Deklination um $3'$ zu groß an. Die obigen Elemente ergeben dagegen eine berechnete Rektaszension von $353^\circ\;14'\;30''$ und eine Deklination von $9^\circ\;12'\;26''$. Jene ist also um $2'\;13''$, diese um $31''$ zu groß. Übrigens ist die beobachtete Deklination ziemlich unsicher, weil der Vergleichsstern sehr ungünstig stand.

N\SourceLeaf{308}och ein paar Beobachtungen aus diesem Monat kann ich Ihnen hier mitteilen.

\par\indent Eintritt von i Virginis am dunklen Mondrand am 4. Juni, $9^{\text{u}}\,3^{\text{m}}\,39^{\text{s}}0$ M. Z., von mir und Professor \textsc{Harding} in derselben Sekunde beobachtet.

\Needspace{8\baselineskip}
\par\indent Eintritt des II$^{\text{ten}}$ Jupitermonds am 22. Juni:


\begin{center}\fontsize{10}{12}\selectfont\setlength{\tabcolsep}{3pt}\renewcommand{\arraystretch}{1.15}\begin{tabular}{ll}
$12^{\text{u}}\,45^{\text{m}}\,9^{\text{s}}$ M. Z. & \textsc{Harding} am $3\frac{1}{4}$-Fuß-\textsc{Dollond}.\\
$12\;45\;32$\quad — & \textsc{Gauss} am 17-Fuß-\textsc{Herschel}.\\
\end{tabular}\end{center}


\par\indent Von Professor \textsc{Hardings} Himmelskarten erwarten wir demnächst das vierte Blatt.

